\documentclass[10pt,a4paper]{amsart}
\usepackage[margin=19mm]{geometry}
\usepackage{amsmath,amssymb,mathtools}
\usepackage[scr=rsfs,cal=euler]{mathalfa}
\usepackage{xfrac}
\usepackage[unicode,hidelinks,bookmarks=false]{hyperref}
\newcommand{\ie}{i.e.\ }
\newcommand{\cf}{cf.\ }
\newcommand{\resp}{respectively\ }
\newcommand{\Subsection}{\subsection}
\providecommand{\nobreakdash}{\nobreak\hbox{-}}
\setlength{\emergencystretch}{2em}
\multlinegap0pt
\usepackage{bm}
\usepackage{bbm}
\usepackage{dsfont}
\usepackage{xspace}
\usepackage{cancel}
\usepackage{mathtools}
\usepackage{amsthm}
\makeatletter
\@ifundefined{theorem}{\newtheorem{theorem}{Theorem}}{}
\@ifundefined{lemma}{\newtheorem{lemma}[theorem]{Lemma}}{}
\@ifundefined{proposition}{\newtheorem{proposition}[theorem]{Proposition}}{}
\makeatother
\newcommand{\Q}{\mathbb{Q}}
\newcommand{\Z}{\mathbb{Z}}
\newcommand{\op}[1]{\operatorname{#1}}
\title[Selmer stability for elliptic curves in Galois $\ell$-extens]{Selmer stability for elliptic curves in Galois $\ell$-extensions}
\author{Siddhi Pathak, Anwesh Ray}
\date{Source: arXiv:2504.15945v2 (CC BY 4.0)}
\begin{document}
\maketitle

\noindent\emph{Source and licence.} Selmer stability for elliptic curves in Galois $\ell$-extensions. Version arXiv:2504.15945v2 (CC BY 4.0), distributed under the Creative Commons Attribution 4.0 International licence, as recorded in the arXiv licence metadata for this exact version and shown on the abstract page https://arxiv.org/abs/2504.15945v2. Full rights notice: licenses/arxiv-2504.15945-CC-BY-4.0.txt.\par\smallskip

\noindent\textit{Editorial scope.} Counted unit: the Proposition whose source label is \texttt{constructing tilde L satisfying SN} in the article (source0.tex lines 621--628, source label \texttt{constructing tilde L satisfying SN}), delivered together with the article's own complete original proof of it (source0.tex lines 629--657, one \texttt{proof} environment of 29 lines). No mathematics is written, repaired or re-derived: statement and proof are the article's own text line for line. Required context retained uncounted: local embedding problem (lines 503--505); local extension is solvable (lines 521--523); step 2 propn (lines 600--602); linear disj lemma (lines 609--615). Every definition, display and label of those lines is retained unchanged. Omitted from this package and not redistributed: 1--502, 506--520, 524--599, 603--608, 616--620, 658--1013. Nothing inside a retained excerpt is omitted or altered: every retained body is delivered exactly as the article's lines are written, so no omission receipt of any kind accompanies this package. Citations: the retained excerpts carry no citation command at all, so no bibliography entry of the article is transcribed and no reference is redistributed. Reference adaptation: none was needed and none was made; every reference inside the delivered unit resolves inside it, so no unresolved reference or citation is produced. Printed slips kept literal: no printed word, symbol, hypothesis or number of the counted statement or of its proof was altered. Layout and class: the article's own class is \texttt{amsart} with packages including amscd, amsmath, amsxtra, amsthm, amssymb, stmaryrd, xr, mathrsfs, mathtools, enumerate, commath, comment, enumitem, xcolor; the wrapper is the lane's pinned \texttt{amsart} editorial wrapper at 10pt on A4, which restates the theorem-like environments the retained text needs and adds the article's own macro definitions that the retained text needs (\texttt{\textbackslash Q}, \texttt{\textbackslash Z}, \texttt{\textbackslash op}), with extra wrapper packages (bm, bbm, dsfont, xspace, cancel, mathtools). The delivered numbering is the unit's own. Assets: no retained span contains a figure environment, a graphics-inclusion command, a PDF-page inclusion, a table or a TikZ picture; no image file is copied into this package, the package declares no asset dependency and no image file is redistributed with it. Licence: version arXiv:2504.15945v2 of this article is distributed under the Creative Commons Attribution 4.0 International licence, as recorded in the arXiv licence metadata for this exact version and shown on the abstract page https://arxiv.org/abs/2504.15945v2. A full rights notice is delivered at licenses/arxiv-2504.15945-CC-BY-4.0.txt. Characters: every retained excerpt is pure ASCII, so the delivered engine, XeLaTeX with its default Latin Modern text font, reports no missing character.
\begin{equation}\label{local embedding problem}
0 \to \mathbb{Z}/\ell\mathbb{Z} \to \widetilde{G}_p \to G_p \to 1
\end{equation}

\begin{proposition}\label{local extension is solvable}
     Suppose \( L/\mathbb{Q} \) is a Galois extension with \(\op{Gal}(L/\Q)\simeq  G \) and canonical surjection \( \varphi : \op{G}_{\Q}\twoheadrightarrow G\) as before. Then, for every prime \( p \), the local embedding problem \eqref{local embedding problem} is solvable. 
\end{proposition}

\begin{proposition}\label{step 2 propn}
    With respect to notation above, there exists $\widetilde{L}/L/\Q$ solving the embedding problem such that $\widetilde{L}$ is ramified at the same set of primes as $L$.
\end{proposition}

\begin{lemma}\label{linear disj lemma}
The following assertions hold:
\begin{enumerate}
    \item[(i)] The fields $F$ and $L$ are linearly disjoint,
    \item[(ii)] $FL$ and $\Q(\mu_\ell, \sqrt[\ell]{p_1}, \dots, \sqrt[\ell]{p_m})$ are linearly disjoint. 
\end{enumerate}
\end{lemma}

\begin{proposition}\label{constructing tilde L satisfying SN}
    There exists $\widetilde{L}/L/\Q$ solving the embedding problem such that:
    \begin{enumerate}
        \item $\widetilde{L}$ satisfies $(\mathfrak{S}_N)$,
        \item $\widetilde{L}$ is ramified at the primes in $\Sigma$ and one additional prime $q$. The prime $q$ can be chosen from a certain set of primes which has positive density.
    \end{enumerate}
     
\end{proposition}

\begin{proof}
From Proposition \ref{step 2 propn}, we obtain a solution \(\widetilde{L}\) to the embedding problem for \(\widetilde{G}\), which is ramified at precisely the same set of primes \(\Sigma = \{p_1, \dots, p_m\}\) as \(L\). Each prime \(p_i\) satisfies \(p_i \equiv 1 \pmod{\ell^N}\), and the decomposition group \(G_{p_i}\) coincides with the inertia group \(\operatorname{I}_{p_i}\) at \(p_i\). Since \(p_i\) is tamely ramified in \(L\), the inertia group \(\operatorname{I}_{p_i}\) is cyclic, and consequently, \(G_{p_i}\) is also cyclic for each \(i = 1, \dots, m\).

Now, fix \(p = p_i\) and let $I_p'$ (resp. $G_p'$) is the inertia (resp. decomposition) group of $p$ in $\widetilde{L}$. On the other hand, we set $\widetilde{I}_p$ (resp. $\widetilde{G}_p$) to be the inverse image of $I_p$ (resp. $G_p$) with respect to the map $\pi: \widetilde{G}\rightarrow G$. Since $I_p=G_p$, it follows that $\widetilde{I}_p=\widetilde{G}_p$. Since the ramification of $p$ is tame, the inertia groups $I_p$ and $I_p'$ are cyclic. Suppose that \(I_p \simeq \mathbb{Z}/\ell^\alpha\mathbb{Z}\). From the proof of Proposition \ref{local extension is solvable}, $\widetilde{G}_p$ is abelian, hence so is $\widetilde{I}_p$. Therefore, there are two possibilities for \(\widetilde{I}_p\).
\begin{description}
    \item[Case 1. \(\widetilde{I}_p \simeq \mathbb{Z}/\ell^\alpha\mathbb{Z} \times \mathbb{Z}/\ell\mathbb{Z}\)] Since $I_p'$ is cyclic and the quotient map from $I_p'$ to $I_p$ is surjective, it follows that $I_p'\simeq I_p$. On the other hand, $G_p'\subseteq \widetilde{G}_p=\widetilde{I}_p$ and thus $G_p'/I_p'$ is either $1$ or $\Z/\ell \Z$. There is an inclusion
    \[G_p'/I_p'\hookrightarrow \widetilde{G}_p/I_p'\xrightarrow{\sim} \Z/\ell \Z\] with respect to which the Frobenius element \(\sigma_p\) maps to some \(c_p \in \mathbb{Z}/\ell\mathbb{Z}\). If \(c_p = 1\), then \(G_p' = I_p'\) and the condition \((\mathfrak{S}_N)\) is satisfied at \(p\). 
    \item[Case 2. \(\widetilde{I}_p \simeq \mathbb{Z}/\ell^{\alpha+1}\mathbb{Z}\)] Since the map $I_p'\rightarrow I_p$ is surjective, it follows that $I_p'=\widetilde{I}_p=\widetilde{G}_p$. This forces $I_p'=G_p'$. In this case, we simply set $c_p:=1$.
\end{description}
If \(c_p = 1\) for all the primes $p$ in Case 1 above, then \(\widetilde{L}\) satisfies property \((\mathfrak{S}_N)\). Otherwise, we must modify \(\widetilde{L}\). Define $X$ as the set of all primes ramified in \(L\) such that $\widetilde{I}_p\simeq I_p\times \Z/\ell \Z$. \\

  Given a prime $q\equiv 1\pmod{\ell^N}$, there exists a unique surjective Galois character \[\chi^{(q)}:\op{Gal}(\Q(\mu_q)/\Q)\rightarrow \Z/\ell\Z.\] Identifying $\op{Gal}(\Q(\mu_q)/\Q)$ with $(\Z/q\Z)^\times$, and the $\ell$-torsion in $(\Z/q\Z)^\times$ with $\Z/\ell\Z$, $\chi^{(q)}$ is given by $\chi^{(q)}(a):=a^{(q-1)/\ell}$. Set \(\mathbb{L} := L\left(\mu_{\ell^N}\cup \{\sqrt[\ell]{p}| p
\in X\}\right)\) and let \(\Delta := \operatorname{Gal}(\mathbb{L}/\mathbb{Q})\). There exists \(\sigma\in \Delta\) such that, for any prime \(q\) unramified in \(\mathbb{L}\) with \(\sigma_q =\sigma\), the following conditions hold:

\begin{enumerate}
    \item \(q \equiv 1 \pmod{\ell^N}\),
    \item \(q\) splits completely in \(L/\mathbb{Q}\),
    \item \(\chi^{(q)}(p) = c_{p}\) for all $p\in X$.
\end{enumerate}
More precisely, $\sigma\in \op{Gal}(\mathbb{L}/\Q)$ is the element defined by:
\begin{itemize}
    \item $\sigma_{|FL}=1$, 
    \item $\sigma(\zeta_\ell)=\zeta_\ell$, 
    \item $\sigma(\sqrt[\ell]{p})=c_p \sqrt[\ell]{p}$ for all primes $p\in X$.
\end{itemize}
Such an element exists since $FL\cap \Q\left(\mu_{\ell^N}\cup \{\sqrt[\ell]{p}| p
\in X\}\right)=\Q$ by Lemma \ref{linear disj lemma}. Let \(\widetilde{\varphi} : \operatorname{G}_{\mathbb{Q}} \to \widetilde{G}\) correspond to \(\widetilde{L}\), and define \(\widetilde{\psi} := \widetilde{\varphi} (\chi^{(q)})^{-1}\), with the corresponding extension denoted \(\widetilde{L}'/L/\mathbb{Q}\). This modified extension solves the embedding problem for \(\widetilde{G}\) while ensuring that property \((\mathfrak{S}_N)\) is satisfied.

\end{proof}

\end{document}
